\Askhsh{24760}{4}
\begin{ylh}{1}
    \item 1.8 Συνέχεια συνάρτησης
    \item 2.6 Συνέπειες του Θεωρήματος Μέσης Τιμής
    \item 2.7 Τοπικά ακρότατα συνάρτησης
    \item 2.8 Κυρτότητα - Σημεία καμπής συνάρτησης.
\end{ylh}
Δίνεται η συνάρτηση $f(x)=e^x-\ln x-\lambda x, x>0$ όπου $\lambda \in \mathbb{R}$. Αν ισχύει $e-\lambda = e^e-1-\lambda e$, να αποδείξετε ότι:
\begin{alist}
\item η $f$ είναι κυρτή. \monades{6}
\item υπάρχει ακριβώς ένα $x_o \in (1,e)$ με $f'(x_o)=0$. \monades{6}
\item για την $f'$ ισχύουν οι υποθέσεις του θεωρήματος Bolzano στο $[1,e]$. \monades{6}
\item η $f$ παρουσιάζει ολικό ακρότατο στο $x_o$ που είναι το $e^{x_o}(1-x_o)+1-\ln x_o$. \monades{7}
\end{alist}